\section{A nonzero example: \texorpdfstring{$\SQS(8)$}{SQS(8)}}\label{sec:app-sqs}

The Steiner quadruple system $\SQS(8)=\AG(3,2)$ is a $3$-$(8,4,1)$ design: its
$14$ blocks are the affine planes (cosets of the seven $2$-dimensional
subspaces) of $\FF_2^{3}$. Its star sequence has $\lambda_1=7$, $\lambda_2=3$,
$\lambda_3=1$, so the star differences are $4$ and $2$ and
$d=\gcd\{4,2\}=2\neq1$; Corollary~\ref{cor:d1} therefore does \emph{not} force
$\Delt=0$. The pointed presentation is $8\times14$. Its Smith form is
$1^5 2^3$, equivalently its $\binom{14}{8}=3003$ maximal minors have gcd $8$.
Together with the exponent-two bound this gives
\[
\Delt(\SQS(8))\cong(\ZZ/2)^{3},\qquad\text{full incidence Smith form } 1^{4}2^{3}4 .
\]
Here $\SQS(8)$ is the smallest nontrivial Steiner quadruple system, and this
nonzero defect shows that the hypothesis $d=1$ of Corollary~\ref{cor:d1} has
genuine content: dropping it admits a nontrivial $2$-group. The verifier
reproduces the Smith forms and modular ranks, rather than enumerating these
minors (Appendix~\ref{sec:app-repro}).
